\subsection{Eigenvalues and eigenvectors}

DEFINITION 1. --- Let E be a vector space over a commutative field K, and u an endomorphism of E. An element x of E is said to be an eigenvector of u if there exists \(\alpha \in K\) such that \(u \cdot x = \alpha x\) ; if \(x \neq 0\) then the scalar \(\alpha\) is called the eigenvalue of u corresponding to the eigenvector x. For every scalar \(\alpha\) , the vector subspace \(V_{\alpha}\) consisting of \(x \in E\) such that \(u \cdot x = \alpha x\) is called the eigenspace of E corresponding to \(\alpha\) .

The geometric multiplicity of the eigenvalue \(\alpha\) is the cardinal dim \(V_{\alpha}\) .

Suppose E is finite dimensional. The eigenvalues of u are those elements \(\alpha\) of K such that the endomorphism \(\alpha \cdot 1 - u\) of E is not injective, in other words (III, p.~524, Prop. 3) such that \(\det(\alpha \cdot 1 - u) = 0\) . But, by the definition of the characteristic polynomial \(\chi_{u}\) of u (III, p.~541, Def. 3), we have \(\det(\alpha \cdot 1 - u) = \chi_{u}(\alpha)\) . Consequently:

PROPOSITION 1. --- Suppose E is finite dimensional. Then an element \(\alpha\) of K is an eigenvalue of the endomorphism \(u\) if and only if it is a root of the characteristic polynomial of \(u\) .

If L is an extension of the field K, then the roots of \(\chi_{u}\) in L are eigenvalues of the endomorphism \(1_{L} \otimes u\) of the L-vector space \(L \otimes_{K} E\) . They are often referred to as eigenvalues of u in L. By abuse of language, we say that all the eigenvalues of u belong to L if this holds for all the eigenvalues of u in an algebraically closed extension of L; this means that \(\chi_{u}\) decomposes into linear factors in L{[}X{]}.

Let \(U\) be a square matrix of order \(n\) with coefficients in \(K\) . Then by definition the characteristic polynomial of \(U\) is

\[
\chi_ {U} (\mathbf {X}) = \det \left(\mathbf {X} \cdot I _ {n} - U\right);
\]

the eigenvalues of U (in an extension L of K) are the roots (in L) of the polynomial \(\chi_{U}\) ; these are also the scalars \(\alpha\) (in L) such that there exists a nonzero solution to the system of linear equations \(UX = \alpha X\) , where X is a column matrix of order n; a column matrix X satisfying this equation is called an eigenvector of U corresponding to \(\alpha\) .

If U is the matrix of an endomorphism u of an n-dimensional vector space with respect to a basis B, then \(\chi_{U} = \chi_{u}\) , the eigenvalues of U are the eigenvalues of u, and the eigenvectors of U are the matrices of the eigenvectors of u with respect to the basis B.

PROPOSITION 2. --- Let u be an endomorphism of a vector space E over a commutative field K; for each scalar \(\alpha\) , let \(V_{\alpha}\) be the eigenspace corresponding to \(\alpha\) . Then the subspaces \(V_{\alpha}\) are closed under u and the sum of the \(V_{\alpha}\) is direct.

The first assertion is clear. By definition the subspace \(V_{\alpha}\) is annihilated by the element X - \(\alpha\) of K{[}X{]}; the X - \(\alpha\) , \(\alpha \in K\) , are irreducible and pairwise non-associate; the second assertion thus follows from VII, p.~8, Th. 1.

